\section{Default style hierarchy}
\label{app:tree}

The default reference tree, hand-built from standard art-history
references and used as the primary evaluation reference throughout
the paper.

\begin{verbatim}
Root
|-- Early_Renaissance
|   |-- Northern_Renaissance
|   `-- High_Renaissance
|       `-- Mannerism_Late_Renaissance
|           `-- Baroque
|               `-- Rococo
|                   `-- Romanticism
|                       |-- Realism
|                       |   |-- Contemporary_Realism
|                       |   `-- Impressionism
|                       |       `-- Post_Impressionism
|                       |           |-- Pointillism
|                       |           |-- Fauvism
|                       |           `-- Cubism
|                       |               |-- Analytical_Cubism
|                       |               `-- Synthetic_Cubism
|                       `-- Symbolism
|                           |-- Art_Nouveau
|                           `-- Expressionism
|                               `-- Abstract_Expressionism
|                                   |-- Action_painting
|                                   |-- Color_Field_Painting
|                                   `-- Minimalism
|-- Pop_Art
|   `-- New_Realism
|-- Ukiyo_e
`-- Naive_Art_Primitivism
\end{verbatim}

\section{Class distribution and tree-distance matrices}
\label{app:data}

\begin{figure}[h]
 \centering
 \includegraphics[width=0.7\linewidth]{Class_Distribtuion.png}
 \caption{Per-style image counts on WikiArt-Refined. Impressionism
 dominates with $13{,}060$ examples; Analytical Cubism, Action
 Painting, and Synthetic Cubism each have fewer than $250$. The
 imbalance ratio between extrema is $133.3\times$.}
 \label{fig:class-distribution}
\end{figure}

\begin{figure}[h]
 \centering
 \includegraphics[width=\linewidth]{dist_mat.png}
 \caption{Tree distance matrices for the three reference hierarchies
 used in this paper: default lineage (left), chronological era
 grouping (centre), flat null (right). Rows and columns are styles in
 the same order across panels.}
 \label{fig:tree-distance-matrices}
\end{figure}

\section{Curvature within the hyperbolic family}
\label{app:phase1}

The body \cref{fig:phase1} shows the dimension scaling for both
geometries; here we add detail on curvature. Sweep configuration:
geometry $\times$ $d \in \{2,4,8,16,32,64\}$ $\times$
$c \in \{0.1, 0.3, 1, 3\}$ for hyperbolic, three seeds. Within the
hyperbolic family, curvature interpolates between the two extremes:
higher $c$ pulls prototypes closer to the boundary of the disk,
where the Poincar\'e metric becomes most curved; this exchanges
$2$--$3$\,pp of top-1 for $4$--$5$\,pp of sibling recall. At $d{=}8$,
the optimal $c$ for top-1 is $0.3$ ($59.8\%$) and the optimal $c$
for sibling recall is $3.0$ ($0.230$). This is consistent with how
hyperbolic capacity is distributed: most volume lies near the
boundary, where the leaves of a tree should live.

\section{Regularizer-strength sweep, full panel set}
\label{app:phase2-full}

\begin{figure}[h]
  \centering
  \includegraphics[width=0.85\linewidth]{phase2_lambda_full.pdf}
  \caption{Full four-panel version of the regularizer-strength sweep
  (body \cref{fig:phase2} drops the third panel for legibility).
  Mean tree distortion decreases monotonically with $\lambda$ in
  both geometries, mirroring the rise in tree-Spearman.}
  \label{fig:phase2-full}
\end{figure}

\section{Training-tree ablation, full panel set}
\label{app:phase3-full}

The body \cref{fig:phase3} summarises the three-tree ablation; the
underlying observation is that training against a tree other than
the evaluation tree provides essentially no useful hierarchical
signal: tree-Spearman against the default tree drops to
$\approx 0.14$ for both non-default training trees, and top-1 climbs
back to the $\lambda{=}0$ baseline because the regularizer no longer
competes with the likelihood.

\begin{figure}[h]
  \centering
  \includegraphics[width=0.85\linewidth]{phase3_tree_variants_full.pdf}
  \caption{Full four-panel version of the training-tree ablation
  (body \cref{fig:phase3} drops the third panel for legibility).
  Tree distortion against the default tree is uniformly higher when
  training against a non-default tree.}
  \label{fig:phase3-full}
\end{figure}

\section{Empirical reference trees and DINOv2 cross-encoder}
\label{app:phase45}

\begin{figure}[h]
  \centering
  \includegraphics[width=0.85\linewidth]{phase4_empirical_tree.pdf}
  \caption{Prototype-distance / tree-distance Spearman against the
  hand-built default tree (blue) and a CLIP-empirical tree (red).
  Both geometries' prototypes align substantially better with the
  empirical tree than with the hand-built one. Caveat: the
  empirical tree is built in Euclidean space from the same CLIP
  features used to train the prototypes, so the alignment with
  prototype distances is partly expected by construction. The
  geometry comparison is the load-bearing claim, not the absolute
  level.}
  \label{fig:phase4}
\end{figure}

\begin{figure}[h]
  \centering
  \includegraphics[width=0.85\linewidth]{phase5_dinov2.pdf}
  \caption{Substituting a DINOv2 ViT-B/14 encoder for CLIP in the
  empirical-tree construction. \textbf{Left:} prototype-tree
  Spearman against three references. The hyperbolic global-Spearman
  lead seen against the CLIP-empirical tree does not replicate
  against the DINOv2-empirical tree: Euclidean prototypes lead at
  every $d \geq 4$. \textbf{Centre:} sibling recall@5 against the
  three reference-tree sibling-set definitions. Hyperbolic leads on
  every reference at every $d \geq 4$. \textbf{Right:}
  pairwise-distance Spearman between reference trees. Hand-built
  correlates only weakly with either encoder-derived tree
  ($0.08$, $0.12$); the two encoder-derived trees correlate $0.50$
  with each other.}
  \label{fig:phase5}
\end{figure}

\paragraph{Tree-vs-tree.} Pairwise-distance Spearman between
hand-built default and CLIP-empirical is $0.08$; between default and
DINOv2-empirical it is $0.12$; between CLIP-empirical and
DINOv2-empirical it is $0.50$. Whatever the hand-built lineage tree
captures, the encoders see only a weak projection of it; whatever
the encoders capture, they capture it similarly. The hand-built
tree is the outlier.

\paragraph{Empirical sibling and cousin recall.} Sibling/cousin sets
used by recall@5 in the body figures are derived from
\texttt{STYLE\_HIERARCHY} (the default tree). To test reference-tree
sensitivity we rebuild sibling and cousin sets from each empirical
tree's binary linkage (siblings: the leaves in the other branch of
the leaf's first merge; cousins: the leaves in the uncle subtree at
the grandparent merge) and rerun the recall computation. The hyperbolic lead is preserved on every reference
tree (sibling recall@5 at $d{=}8$, default / CLIP-empirical /
DINOv2-empirical: Eu $0.149 / 0.356 / 0.322$, Hy $0.195 / 0.416 /
0.380$).

\paragraph{Global-Spearman direction is encoder-specific.} Against
the DINOv2-empirical tree, Euclidean prototypes lead at every
$d \geq 4$, with $\rho_{\textrm{Eu}} - \rho_{\textrm{Hy}}$ growing to
$+0.06$ at $d{=}64$ ($0.413$ vs.\ $0.357$). This is a negative result
for the strongest reading of the CLIP-empirical finding and
motivates the global-fidelity caveat in the body.

\section{Class-imbalance robustness}
\label{app:imbalance}

WikiArt is heavily skewed (Impressionism has $133\times$ more
training examples than Action Painting). The default sweep uses
unweighted cross-entropy; rerunning the $d{=}8$ winners of each
geometry with inverse-frequency class-weighted cross-entropy (three
seeds) shifts absolute numbers as expected (top-1 drops by roughly
$6$\,pp for both: Eu $64.3 \to 58.3\%$, Hy $59.8 \to 53.4\%$; while
balanced accuracy rises sharply: Eu $56.4 \to 65.6\%$, Hy
$44.5 \to 62.1\%$), but the geometry gap is preserved on every axis. Top-1 remains $5$\,pp Euclidean, sibling recall remains
$4$\,pp hyperbolic, and class-center / tree Spearman shifts by less
than a hundredth. The balanced-accuracy gap narrows from $11.9$ to
$3.5$\,pp, suggesting hyperbolic's main classification weakness in
the default setup was disproportionately on rare classes; it remains
in Euclidean's favour.

\section{Confusion structure}
\label{app:confusion}

\begin{figure}[h]
  \centering
  \includegraphics[width=0.85\linewidth]{confusion_block_d8.pdf}
  \caption{Row-normalized confusion matrices at $d{=}8$, with rows
  and columns permuted by a depth-first traversal of the default
  style tree so that hierarchically adjacent styles sit next to each
  other on both axes. Both models concentrate mistakes in
  near-diagonal blocks: errors predominantly fall on tree-adjacent
  styles. The hyperbolic block structure is visibly
  softer near the diagonal: more confusion mass spreads into
  immediate siblings, less mass jumps to distant styles. This is
  the same fact the sibling-recall numbers report, in a form a
  viewer can absorb at a glance.}
  \label{fig:confusion}
\end{figure}

\section{Reproducibility}
\label{app:repro}

The full sweep CSV records, for every config, the training
hyperparameters and every evaluation metric (150 rows). To reproduce:
\begin{verbatim}
python scripts/sweep.py --phase 1 --device mps  # ~14 min, 90 configs
python scripts/sweep.py --phase 2 --device mps  # ~10 min, 48 configs
python scripts/sweep.py --phase 3 --device mps  # ~3 min,  18 configs
python scripts/baselines.py                     # logreg, kNN baselines
python scripts/significance.py                  # paired-t, sign tests
python scripts/make_figures.py                  # regenerate figures
\end{verbatim}
Code, sweep configurations, and exact metric implementations will be
released after de-anonymisation.
